% AUTO-GENERATED by writeup/make_rehaul_disclosures.py -- do not edit; re-run the script.
% \rhds... macros: the ends of the feeds and the data fixes of the final pooled design (Section 3, note E28).
% Run files: results/window_sweep_2014start_2026-10-22/w3000; scored bars from 2014-01-02.
\providecommand{\rhdsSampleEnd}{}\renewcommand{\rhdsSampleEnd}{30 April 2024}  % last day of the data (the end of the VIX splice, loading.VIX_SPLICE_END)
\providecommand{\rhdsVixVendorLast}{}\renewcommand{\rhdsVixVendorLast}{12 February 2024}  % last vendor print of VIX, VIX3M and VVIX (data/vix_and_voldemand.parquet)
\providecommand{\rhdsVixSpliceDays}{}\renewcommand{\rhdsVixSpliceDays}{54}  % trading days after it to the end of the data (16:00 futures print), spliced from Yahoo
\providecommand{\rhdsVixOverlapStart}{}\renewcommand{\rhdsVixOverlapStart}{December 2023}  % first month of the stamps both the vendor file and the Yahoo hourly feed carry
\providecommand{\rhdsVixAgree}{}\renewcommand{\rhdsVixAgree}{96}  % percent of overlap stamps on which Yahoo and the vendor print the same VIX (to the second decimal)
\providecommand{\rhdsVixThreeMAgree}{}\renewcommand{\rhdsVixThreeMAgree}{88}  % the same, VIX3M
\providecommand{\rhdsVvixAgree}{}\renewcommand{\rhdsVvixAgree}{100}  % the same, VVIX
\providecommand{\rhdsVdLast}{}\renewcommand{\rhdsVdLast}{31 August 2023}  % last print of the four volatility-demand series
\providecommand{\rhdsFfillCap}{}\renewcommand{\rhdsFfillCap}{26}  % forward-fill cap of the exogenous series, intervals (executor.FFILL_LIMIT under APPLY_FFILL_LIMIT)
\providecommand{\rhdsFomcFileLast}{}\renewcommand{\rhdsFomcFileLast}{1 November 2023}  % last scheduled FOMC statement in data/releases.parquet
\providecommand{\rhdsFomcAdded}{}\renewcommand{\rhdsFomcAdded}{three}  % scheduled statements added after it (loading.FOMC_MISSING_STATEMENTS), as a word
\providecommand{\rhdsInvalidPrints}{}\renewcommand{\rhdsInvalidPrints}{21{,}943}  % prints set to missing by the domain rule (loading.null_invalid_prints)
\providecommand{\rhdsInvalidVolume}{}\renewcommand{\rhdsInvalidVolume}{21{,}852}  % ... of them zero futures volumes (sumvolume)
\providecommand{\rhdsValueCols}{}\renewcommand{\rhdsValueCols}{276}  % value-average columns adj_<c>_ma_<k> of the full predictor set (final design)
\providecommand{\rhdsUnobsCols}{}\renewcommand{\rhdsUnobsCols}{134}  % of them set to 0 by the no-observation rule on some row of the design (run_info.json design log)
\providecommand{\rhdsUnobsColsEval}{}\renewcommand{\rhdsUnobsColsEval}{113}  % of them set to 0 on some scored bar from the first evaluation day (recomputed; gate: design log)
\providecommand{\rhdsEvalStart}{}\renewcommand{\rhdsEvalStart}{2 January 2014}  % first evaluation day of the final runs: the first scored bar's day (final_run.FIRST_DAY)


\section{Data}
\label{sec:rh_data}
\begin{itemize}
\item \textbf{E-mini futures and realized variance.} We obtain one-minute prices of the front-month E-mini S\&P~500
  futures contract (the most actively traded contract), which trades nearly 24 hours a day, from 1998 through
  \rhdsSampleEnd, and aggregate them on a grid of half-hour intervals. Trading days are indexed by $t = 1, \dots, T$
  and the half-hour intervals of day $t$ by $i = 1, \dots, M$; every interval is stamped by its end in Eastern time,
  and $i = c$ denotes the closing interval, 15:30--16:00~ET. With $r_{t,i,k}$ the one-minute log-price change of the
  front-month contract in the $k$-th non-missing minute of interval $(t,i)$, $k = 1, \dots, N_{t,i}$
  ($N_{t,i} \le 30$), the realized variance of the interval is \citep{AndersenBollerslevDieboldLabys2003}
  \[
    \RV_{t,i} = \sum_{k=1}^{N_{t,i}} r_{t,i,k}^2 ,
  \]
  which we forecast on every regular-hours interval, the bars ending 10:00 to 16:00~ET
  (Section~\ref{sec:rh_methods_setup}).
\item \textbf{Exogenous series.} The models use \rhdmNSeriesPooled{} exogenous series in \bkNBuckets{} groups.
  Appendix Table~\ref{tab:rh_dm_dictionary} defines each series, and Appendix~\ref{app:rh_dp_sources} gives the
  sources and the construction.
  \begin{itemize}
  \item \textbf{Realized moments} (\bkSizeMom{} series): one-minute sums of the futures over the interval other than
    $\RV_{t,i}$, the target: the sums of the returns, of their absolute values, cubes and fourth powers, the positive
    realized semivariance, the bipower variation and the first-order realized autocovariance.
  \item \textbf{Liquidity} (\bkSizeLiq{} series): the futures volume, the number of observed minutes
    $N_{t,i}$\ifrawseries{} (\sername{numobs})\fi, and equal- and value-weighted cross-sectional averages over the
    S\&P~500 constituents of turnover (total, buyer- and seller-initiated), the effective and quoted spreads, and
    order-flow imbalance.
  \item \textbf{Equal- and value-weighted market moments} (\bkSizeMew{} and \bkSizeMvw{} series): the same averages of
    realized moments of the constituents' 5-minute returns.
  \item \textbf{Sentiment} (\bkSizeSent{} series): StockTwits attention (posts in the half hour), the sentiment score of
    the posts tagged bullish or bearish, and their number.
  \item \textbf{Implied volatility} (\bkSizeIv{} series): the Cboe VIX, VIX3M and VVIX indices.
  \item \textbf{Volatility demand} (\bkSizeVd{} series): the net vega bought by customers other than market makers in
    Cboe option volume.
  \item \textbf{FOMC} (\bkSizeFomc{} series): an indicator for a scheduled Federal Open Market Committee statement on
    day $t$ and the proximity measures $1/(1+n^{+}_{t,i})$ and $1/(1+n^{-}_{t,i})$, with $n^{+}_{t,i}$ and
    $n^{-}_{t,i}$ the number of intervals until the next and since the last scheduled statement. The proximity
    measures follow the pre-announced meeting calendar, so all three series are known before the interval begins;
    the meeting scheduled for \rhdsFomcSchedMeeting{} is the next statement until the action announced on
    \rhdsFomcReplaced{} replaced it (Appendix~\ref{app:rh_dp_fomc}).
  \end{itemize}
  Our data end on \rhdsSampleEnd. The file of the implied volatility indices ends on \rhdsVixVendorLast{} and is
  extended with the hourly closes of a free feed; the FOMC release file ends with the statement of \rhdsFomcFileLast{}
  and is completed with the \rhdsFomcAdded{} scheduled statements that follow; the volatility-demand feed ends on
  \rhdsVdLast{} (Appendix~\ref{app:rh_dp_sources}).
\item \textbf{Interval grid, sessions and cleaning.}
  \begin{itemize}
  \item The interval grid $\mathcal{G}$ holds the intervals with an observed realized variance outside the weekend
    break, Friday 20:00 to Sunday 18:30. Lag operators step back through $\mathcal{G}$, overnight sessions included:
    $(t, i - j)$ is the $j$-th interval of $\mathcal{G}$ before $(t, i)$, which may lie on an earlier day.
  \item Early-close sessions stay in $\mathcal{G}$ but leave the estimation and evaluation samples of the
    models and the option-trading sample; their intervals after the 13:00~ET close also leave the window of the
    target's diurnal factor (Section~\ref{sec:rh_methods_setup}).
  \item On NYSE holidays the E-mini trades a shortened session that ends at about 13:00~ET. The regular-hours intervals
    of these sessions stay in $\mathcal{G}$, in the \rhorNBars-bar estimation sample and in the window of the target's
    diurnal factor, but are not scored; dropping them from the sample and the window changes the every-bar
    comparisons of the full predictor set with the HAR benchmark by \rhholEffCABarPpRng{} percentage points in the
    linear models (Appendix~\ref{app:rh_robust_holiday}).
  \item A print outside the domain of its series, such as a nonpositive spread or futures volume, is set to missing.
    For each exogenous series $u$ we record the availability indicator $a_{t,i} = \mathbf{1}\{u \text{ observed in }
    (t,i)\}$ before any filling, forward-fill $u$ over at most \rhdsFfillCap{} intervals to compute its trailing
    statistics, and impute every interval with $a_{t,i} = 0$ (step~4 of Section~\ref{sec:rh_predictors}). A trailing
    average that holds no observed interval enters at its neutral value (step~6), as the averages of the
    volatility-demand series do after their feed stops (Appendix~\ref{app:rh_dp_sessions}).
  \end{itemize}
\item \textbf{Samples.} Appendix Table~\ref{tab:rh_dm_samples} lists the samples. The estimation samples and the
  evaluation span are set out in Section~\ref{sec:rh_exp_setup}, the option-trading sample in
  Section~\ref{sec:rh_econ_strategy}, and the sample of the closing-interval models in
  Appendix~\ref{app:rh_closing_sample}.
\end{itemize}

